\documentclass{article}
\usepackage[margin=1in]{geometry}
\usepackage{amsmath,amssymb,amsthm}
\usepackage{booktabs,longtable}
\usepackage{hyperref}

\newtheorem{definition}{Definition}
\newtheorem{theorem}[definition]{Theorem}
\newtheorem{lemma}[definition]{Lemma}
\newtheorem{proposition}[definition]{Proposition}
\newtheorem{convention}[definition]{Convention}

\newcommand{\source}[1]{\par\noindent\textbf{Source.} #1}
\newcommand{\lean}[1]{\par\noindent\textbf{Lean target.} \texttt{\detokenize{#1}}}
\newcommand{\status}[1]{\par\noindent\textbf{Status.} \texttt{\detokenize{#1}}}
\newcommand{\depends}[1]{\par\noindent\textbf{Depends on.} #1}
\newcommand{\cA}{\mathcal A}
\newcommand{\Ftwo}{\mathbb F_2}
\newcommand{\tr}{\operatorname{tr}}

\begin{document}

\section*{Blueprint 007: Finite Pauli Algebra and Quasi-Local Interface}

\subsection*{Scope}

This note promotes targets \texttt{PAULI-002} and \texttt{QL-001} from the
master dependency blueprint to a dedicated prerequisite informal blueprint
note.  It records the finite Pauli matrix basis, the normalized trace
coefficient formula, and the quasi-local local-net interface needed before
formalizing positivity, GNS data, represented region algebras, or type
classification.

This note does not introduce a new stabilizer-state positivity theorem.  It
only fixes the finite Pauli algebra and quasi-local algebra carriers on which
later canonical states will live.

\subsection*{Conventions}

\begin{convention}[Index and support conventions]
Let \(\Gamma\) be a finite or countably infinite qubit index set.  For a finite
subset \(F\subset \Gamma\), write
\[
  V_F=(F\to\Ftwo)\oplus(F\to\Ftwo).
\]
If \(F\subset G\), identify \(V_F\) with the subspace of \(V_G\) of labels
that vanish on \(G\setminus F\).  The symplectic form is the finite-support
form already recorded in \texttt{SYS-001} and \texttt{PAULI-001}.
\source{Gottesman, \emph{Surviving as a Quantum Computer in a Classical
World}, Chapter 3.5, Definitions 3.11 and 3.13, PDF pp.~55--56; project proof
note
\path{../runs/2026-06-18-stabilizer-algebra-classification/artifacts/004-internal-stabilizer-reduction-type-classification-R1.tex},
``Region Pauli label space''.}
\status{Convention for \texttt{PAULI-002} and \texttt{QL-001}.}
\end{convention}

\begin{convention}[Finite Pauli representatives]
On one qubit use the matrices
\[
 I=\begin{pmatrix}1&0\\0&1\end{pmatrix},\qquad
 X=\begin{pmatrix}0&1\\1&0\end{pmatrix},\qquad
 Z=\begin{pmatrix}1&0\\0&-1\end{pmatrix}.
\]
For \(v=(x,z)\in V_F\), define the finite Pauli representative
\[
  P_F(v)=\bigotimes_{j\in F} X^{x_j} Z^{z_j}
  \in \bigotimes_{j\in F}M_2(\mathbb C).
\]
This section fixes phases by the ordered product \(X^{x_j}Z^{z_j}\).  The
commutation relation is
\[
  P_F(v)P_F(w)=(-1)^{\sigma(v,w)}P_F(w)P_F(v),
\]
where \(\sigma\) is the binary symplectic form.  Multiplication may acquire
the usual section-dependent sign in the label \(v+w\); later stabilizer
coefficient rules depend only on the fixed section and the commutation form.
\source{Gottesman, Chapter 3.5, Definitions 3.11 and 3.13, PDF pp.~55--56,
for the binary symplectic Pauli convention; project proof note 004, Section 2,
for the finite-support region Pauli convention.}
\status{Phase convention for this blueprint.}
\end{convention}

\begin{definition}[FiniteLocalAlgebra, supporting ID: PAULI-002, project proof note 004, region Pauli label space and product form; standard finite matrix trace calculus]
For a finite \(F\subset\Gamma\), define
\[
  \cA_F=\bigotimes_{j\in F} M_2(\mathbb C),
\]
with the convention \(\cA_\emptyset=\mathbb C\).  Its normalized trace is
\[
  \tr_F=\bigotimes_{j\in F}\tr_2,\qquad
  \tr_2(a)=\frac{1}{2}\operatorname{Tr}(a).
\]
\source{Project proof note 004, ``Region Pauli label space'' and ``Product
form after symplectic basis reduction''; standard finite matrix trace
calculus, to be backed in Lean by mathlib matrix traces and finite tensor
products.}
\lean{Formalization.Stabilizer.FinitePauliAlgebra.LocalAlgebra}
\status{Definition supporting \texttt{PAULI-002}.}
\end{definition}

\begin{lemma}[PauliBasisTraceOrthogonality, ID: PAULI-002, project proof note 004 Sections 2--4; Gottesman Chapter 3.5, PDF pp.~55--56]
For every finite \(F\subset\Gamma\), the family
\[
  \{P_F(v):v\in V_F\}
\]
is a complex linear basis of \(\cA_F\).  It is orthonormal for the normalized
Hilbert--Schmidt inner product
\[
  \langle a,b\rangle_F=\tr_F(a^*b),
\]
so every \(a\in\cA_F\) has the Pauli expansion
\[
  a=\sum_{v\in V_F}\tr_F(P_F(v)^*a)\,P_F(v).
\]
Equivalently, the Pauli coefficient of \(a\) at \(v\) is
\[
  \widehat a(v)=\tr_F(P_F(v)^*a).
\]
\source{Finite-dimensional Pauli calculus used in project proof note 004,
Sections 2--4, especially the local coefficient formula and product-form
argument; Gottesman, Chapter 3.5, PDF pp.~55--56, for Pauli labels and
binary symplectic commutation.}
\depends{\texttt{SYS-001}, \texttt{PAULI-001}.}
\lean{Formalization.Stabilizer.FinitePauliAlgebra.pauliBasis}
\status{blueprint_note_complete}
\end{lemma}

\begin{proof}
For one qubit, the four matrices \(I,X,Z,XZ\) are linearly independent in
\(M_2(\mathbb C)\).  Since \(M_2(\mathbb C)\) has complex dimension \(4\), they
form a basis.  A direct trace calculation gives
\[
  \tr_2(P(v)^*P(w))=\delta_{v,w}
\]
for \(v,w\in\Ftwo^2\).

For a finite set \(F\), the tensor products of the one-site basis elements are
indexed by \(V_F\), and their number is \(4^{|F|}\), the dimension of
\(\bigotimes_{j\in F}M_2(\mathbb C)\).  Tensor products of one-site bases form
a basis of the finite tensor product algebra.  The normalized trace is
multiplicative on elementary tensors, hence
\[
  \tr_F(P_F(v)^*P_F(w))
  =
  \prod_{j\in F}\tr_2(P(v_j)^*P(w_j))
  =
  \delta_{v,w}.
\]
The coefficient formula is the usual orthonormal-basis expansion with respect
to \(\langle a,b\rangle_F=\tr_F(a^*b)\).
\end{proof}

\begin{lemma}[TraceCoeffInclude, supporting ID: PAULI-002, project proof note 004, crossing stabilizers do not change local coefficients; finite tensor trace calculus]
Let \(F\subset G\subset\Gamma\) be finite, and let
\[
  \iota_{F,G}:\cA_F\to\cA_G
\]
be the unital inclusion \(a\mapsto a\otimes 1_{G\setminus F}\).  If
\(v\in V_F\), then
\[
  \tr_G(P_G(v)^*\iota_{F,G}(a))=\tr_F(P_F(v)^*a).
\]
If \(w\in V_G\) has nonzero support in \(G\setminus F\), then
\[
  \tr_G(P_G(w)^*\iota_{F,G}(a))=0.
\]
\source{Project proof note 004, ``Crossing stabilizers do not change the local
coefficients'', for the coefficient-level locality principle; finite tensor
trace calculus as in \texttt{PAULI-002}.}
\depends{\texttt{PAULI-002}.}
\lean{Formalization.Stabilizer.FinitePauliAlgebra.traceCoeff_include}
\status{blueprint_note_complete}
\end{lemma}

\begin{proof}
Write \(\iota_{F,G}(a)=a\otimes 1_{G\setminus F}\).  For \(v\in V_F\), the
Pauli representative factors as \(P_G(v)=P_F(v)\otimes 1\), so multiplicativity
and \(\tr_{G\setminus F}(1)=1\) give the first identity.  If \(w\) has a
nontrivial one-site Pauli factor outside \(F\), then the trace over that
outside site is zero, and the entire tensor trace is zero.
\end{proof}

\begin{lemma}[FiniteLocalInclusionAPI, supporting ID: QL-001, project proof note 004, region Pauli label space; finite tensor trace calculus]
Let \(F\subseteq G\subseteq K\subseteq\Gamma\) be finite.  The finite
inclusion
\[
  \iota_{F,G}:\cA_F\to\cA_G,\qquad
  a\mapsto a\otimes 1_{G\setminus F},
\]
is an injective unital \(*\)-homomorphism, and
\[
  \iota_{F,F}=\operatorname{id}_{\cA_F},\qquad
  \iota_{G,K}\circ\iota_{F,G}=\iota_{F,K}.
\]
\source{Project proof note 004, ``Region Pauli label space''; finite tensor
trace calculus as in \texttt{PAULI-002}.}
\depends{\texttt{PAULI-002}.}
\lean{Formalization.Stabilizer.QuasiLocal.finiteInclusionAPI}
\status{blueprint_note_complete}
\end{lemma}

\begin{proof}
The map \(\iota_{F,G}\) tensors every elementary tensor in \(\cA_F\) with the
identity over the new tensor factors \(G\setminus F\), so it preserves
products, adjoints, scalar multiplication, and the unit.  Injectivity follows
from the normalized partial trace over the added factors: the map
\[
  E_{G,F}=\operatorname{id}_{\cA_F}\otimes \tr_{G\setminus F}:\cA_G\to\cA_F
\]
satisfies \(E_{G,F}(\iota_{F,G}(a))=a\).  The identity and compatibility
equations are the associativity of finite tensor products together with
tensoring by identities.
\end{proof}

\begin{definition}[QubitLocalNetInterface, ID: QL-001, Blackadar III.3.1.1--III.3.1.3, PDF pp.~301--302]
The Lean-facing quasi-local object used in this project is the following
project-local local-net interface.  It packages:
\begin{enumerate}
\item the finite local \(C^*\)-algebras
\[
  \cA_F=\bigotimes_{j\in F}M_2(\mathbb C)
\]
for finite \(F\subseteq\Gamma\);
\item the injective unital \(*\)-homomorphisms
\(\iota_{F,G}:\cA_F\to\cA_G\) for \(F\subseteq G\), with
\(\iota_{F,F}=\operatorname{id}\) and
\(\iota_{G,K}\circ\iota_{F,G}=\iota_{F,K}\);
\item a global unital \(C^*\)-algebra \(\cA_\Gamma\) and injective unital
\(*\)-homomorphisms \(j_F:\cA_F\to\cA_\Gamma\) satisfying
\[
  j_G\circ\iota_{F,G}=j_F;
\]
\item the algebraic local union
\[
  \cA_{\Gamma}^{\mathrm{loc}}
  =
  \bigcup_{F\Subset\Gamma} j_F(\cA_F)
  \subseteq \cA_\Gamma,
\]
which is a unital \(*\)-subalgebra whose norm closure is all of
\(\cA_\Gamma\);
\item for every region \(\Lambda\subseteq\Gamma\), a region \(C^*\)-subalgebra
\[
  \cA_\Lambda
  =
  \overline{\bigcup_{F\Subset\Lambda} j_F(\cA_F)}^{\|\cdot\|}
  \leq \cA_\Gamma,
\]
with its inclusion into the global carrier \(\cA_\Gamma\).
\end{enumerate}
Thus the full \(C^*\)-inductive-limit carrier may be constructed later, but
all Lean targets in the state, GNS, and represented-region layers use these
named fields and not an unspecified carrier.
\source{Project proof note 004, ``Region Pauli label space'', for the local
quasi-local algebra as the norm closure of finite tensor products; Blackadar,
\emph{Operator Algebras}, III.3.1.1--III.3.1.3, PDF pp.~301--302, for the
infinite tensor product/product-state operator-algebra setting later consumed
by GNS and product-factor targets.}
\depends{\texttt{PAULI-002}.}
\lean{Formalization.Stabilizer.QuasiLocal.QubitLocalNet}
\status{blueprint_note_complete}
\end{definition}

\begin{proof}[Informal well-formedness check]
Finite subsets of \(\Gamma\) are directed by union.  The finite inclusion API
above gives a directed system of unital \(C^*\)-algebras with injective
connecting maps.  The Blackadar/Takesaki inductive-limit construction supplies
the global carrier \(\cA_\Gamma\), the canonical maps \(j_F\), and the density
of the algebraic local union.  For each \(\Lambda\), the displayed closure is
a \(C^*\)-subalgebra of \(\cA_\Gamma\) because it is the norm closure of a
unital \(*\)-subalgebra of \(\cA_\Gamma\).  These data are exactly the fields
listed in the interface.
\end{proof}

\begin{lemma}[RegionSubalgebraIsotony, supporting ID: QL-001, project proof note 004 Section 2; Blackadar III.3.1.1--III.3.1.3, PDF pp.~301--302]
If \(\Lambda\subseteq\Lambda'\subseteq\Gamma\), then the region subalgebras in
the \(\texttt{QL-001}\) interface satisfy
\[
  \cA_\Lambda\leq \cA_{\Lambda'}\leq \cA_\Gamma.
\]
Equivalently, the inclusion \(\cA_\Lambda\hookrightarrow\cA_\Gamma\) factors
through the inclusion \(\cA_{\Lambda'}\hookrightarrow\cA_\Gamma\).
\source{Project proof note 004, Section 2, for region-local finite support;
Blackadar III.3.1.1--III.3.1.3, PDF pp.~301--302, for the quasi-local
closure model.}
\depends{\texttt{QL-001}.}
\lean{Formalization.Stabilizer.QuasiLocal.region_isotone}
\status{blueprint_note_complete}
\end{lemma}

\begin{proof}
Every finite \(F\subseteq\Lambda\) is also a finite subset of
\(\Lambda'\).  Therefore the local generating union for \(\cA_\Lambda\) is a
subset of the local generating union for \(\cA_{\Lambda'}\).  Taking norm
closures inside the same global carrier preserves inclusion, giving the
displayed isotony.
\end{proof}

\begin{lemma}[DenseLocalUnion, supporting ID: QL-001, Blackadar III.3.1.1--III.3.1.3, PDF pp.~301--302; project proof note 004 Section 2]
The local union
\[
  \cA_{\Gamma}^{\mathrm{loc}}
  =
  \bigcup_{F\Subset\Gamma} j_F(\cA_F)
\]
is norm dense in the global carrier \(\cA_\Gamma\).  More generally, for a
region \(\Lambda\), the union
\[
  \bigcup_{F\Subset\Lambda} j_F(\cA_F)
\]
is norm dense in \(\cA_\Lambda\).
\source{Blackadar III.3.1.1--III.3.1.3, PDF pp.~301--302; project proof note
004, Section 2.}
\depends{\texttt{QL-001}.}
\lean{Formalization.Stabilizer.QuasiLocal.dense_localUnion}
\status{blueprint_note_complete}
\end{lemma}

\begin{proof}
Global density is a field of the \(\texttt{QL-001}\) interface and is the
completion property of the inductive limit.  The regional statement is the
definition of \(\cA_\Lambda\) as the norm closure of the corresponding local
generating union inside \(\cA_\Gamma\).
\end{proof}

\begin{theorem}[CompatibleLocalStateExtension, supporting ID: QL-001, Takesaki III Definition XIV.1.3, PDF p.~105, Takesaki I Proposition I.9.12, PDF p.~59]
Let \((\varphi_F)_{F\Subset\Gamma}\) be a family such that each
\(\varphi_F:\cA_F\to\mathbb C\) is a positive unital linear functional and,
whenever \(F\subseteq G\),
\[
  \varphi_G(\iota_{F,G}(a))=\varphi_F(a),\qquad a\in\cA_F.
\]
Then there is a unique state
\[
  \varphi:\cA_\Gamma\to\mathbb C
\]
such that
\[
  \varphi(j_F(a))=\varphi_F(a)
\]
for every finite \(F\subseteq\Gamma\) and \(a\in\cA_F\).  The same uniqueness
and existence statement holds on each region algebra: if
\(\Lambda\subseteq\Gamma\), then the compatible subfamily with
\(F\subseteq\Lambda\) extends to a unique state
\(\varphi_\Lambda:\cA_\Lambda\to\mathbb C\).
\source{Takesaki III, Definition XIV.1.3, PDF p.~105, for the
\(C^*\)-inductive-limit completion; Takesaki I, Proposition I.9.12, PDF
p.~59, for continuity of positive functionals; Blackadar III.3.1.1--III.3.1.3,
PDF pp.~301--302, for the quasi-local tensor-product model.}
\depends{\texttt{QL-001}.}
\lean{Formalization.Stabilizer.QuasiLocal.compatibleState_extension_unique}
\status{blueprint_note_complete}
\end{theorem}

\begin{proof}
Define \(\varphi_0\) on the algebraic local union by
\(\varphi_0(j_F(a))=\varphi_F(a)\).  This is well defined because if the same
local element is represented in two finite regions \(F\) and \(G\), then it is
also represented in \(F\cup G\), and compatibility with the two inclusions into
\(F\cup G\) gives the same scalar.  Linearity is checked after moving two
summands into a common finite region.

If \(x=b^*b\) lies in the local union, choose a finite \(F\) containing a
representative of \(b\).  Then
\(\varphi_0(x)=\varphi_F(b_F^*b_F)\geq0\), so \(\varphi_0\) is positive on the
local \(*\)-algebra, and \(\varphi_0(1)=1\).  Positivity and unitality imply
\(|\varphi_0(x)|\leq\|x\|\) on the local union by the standard positive
functional norm estimate.  Hence \(\varphi_0\) extends uniquely by continuity
to the norm closure, which is \(\cA_\Gamma\) by \(\texttt{DenseLocalUnion}\).
The extension is still positive because the positive cone is norm closed and
is still unital because the unit is local.  Any other state with the same
finite restrictions agrees on the dense local union, hence agrees everywhere.
The proof for a region \(\Lambda\) is identical using the dense local union
inside \(\cA_\Lambda\).
\end{proof}

\subsection*{Lean Extraction Notes}

\begin{longtable}{@{}p{0.18\linewidth}p{0.34\linewidth}p{0.38\linewidth}@{}}
\toprule
Target & Intended file & Notes \\
\midrule
\endhead
\texttt{PAULI-002} &
\texttt{Formalization/Stabilizer/PauliAlgebra.lean} &
Start with finite matrices and normalized trace coefficients.  Avoid infinite
systems until the finite basis and inclusion lemmas are stable. \\
\texttt{QL-001} &
\texttt{Formalization/Stabilizer/QuasiLocal.lean} &
Use the explicit \texttt{QubitLocalNet} interface above: finite algebras,
injective inclusions, global carrier, region subalgebras, isotony, dense local
union, and compatible-state extension. \\
\bottomrule
\end{longtable}

\end{document}
